\scanpage{30}%
\mspage{15}%
\[
y^2=x^3-x \qquad (\text{si }c_4=c_6=0) \tag*{(6.7)}
\]

c) $p=2$ : ici, on met la courbe sour la forme (6.3) et on résoud (6.4) avec $a_3=1$, $a_4=a_6=0$ pour obtenir la forme type
\[
y^2+y=x^3 \qquad (\text{si }c_4=c_6=0) \, . \tag*{(6.8)}
\]

Pour compléter le tableau, reste à décrire les automorphismes des courbes elliptiques.

\uline{Proposition} 6.9. \quad \uline{Soit} $E$ \uline{une courbe de genre un sur une base} $S$, \uline{sans fibre additive}.

(I) \quad \uline{Si} $c_4$ et $c_6$ \uline{sont inversibles,} \uline{les seuls automorphismes}\margin{\mbox{\begingroup\edef\encoutersep{\the\dimexpr.5\pgflinewidth\relax}\setbox0=\hbox{\shortstack[l]{\handwritten{faux si $S$}\\\handwritten{\unsure{pas} connexe}}}\edef\encwda{\the\wd0}\edef\enchta{\the\ht0}\edef\encdpa{\the\dp0}\rlap{\smash{\tikz[baseline=0pt]{\begin{scope}[shift={(\the\dimexpr \encwda/2\relax,\the\dimexpr -\encdpa/2+\enchta/2\relax)}]\coordinate (enc-rem-north-east) at (\the\dimexpr \encoutersep+\encwda/2+2pt\relax,\the\dimexpr \encdpa/2+\enchta/2+\encoutersep+2pt\relax);\coordinate (enc-rem-north-west) at (\the\dimexpr -\encoutersep-\encwda/2-2pt\relax,\the\dimexpr \encdpa/2+\enchta/2+\encoutersep+2pt\relax);\coordinate (enc-rem-south-east) at (\the\dimexpr \encoutersep+\encwda/2+2pt\relax,\the\dimexpr -\encdpa/2-\enchta/2-\encoutersep-2pt\relax);\coordinate (enc-rem-south-west) at (\the\dimexpr -\encoutersep-\encwda/2-2pt\relax,\the\dimexpr -\encdpa/2-\enchta/2-\encoutersep-2pt\relax);\coordinate (enc-rem-west) at (\the\dimexpr -\encoutersep-\encwda/2-2pt\relax,0pt);\handwritten{\draw[thin,gray,rounded corners=4pt] (enc-rem-north-east) -- (enc-rem-south-east) -- (enc-rem-south-west) -- (enc-rem-north-west) -- ([xshift=-0.4em]enc-rem-north-east);\draw[thin] ([xshift=0.35em,yshift=0.5em]enc-rem-north-west) -- ([xshift=0.2em]enc-rem-west);}\end{scope}\pgfresetboundingbox\path[use as bounding box] (0pt,\the\dimexpr-\encdpa\relax) rectangle (\encwda,\enchta);}}}\copy0\endgroup}} \uline{de} $E$ \uline{sont l'identité et la symétrie}.

(II) \quad \uline{Si} $2$ et $3$ \uline{sont inversibles}, \uline{l'application canonique}
\diagram{\[
\vcenter{\hbox{\begin{tikzcd}[ampersand replacement=\&,column sep=1.9em]
 \text{\uline{Aut}}(E) \handwritten{\arrow[r]} \& G_m
\end{tikzcd}}}\tag*{(6.6.)}
\]}
\uline{donnée par l'action sur les différentielles invariantes}, \uline{est une immersion fermée}.

(III) \uline{Si} $2$ \uline{et} $3$ \uline{sont inversibles}, \uline{et} $c_4=0$ (resp. $c_6=0$) \uline{l'image de} (6.6) \uline{est} $\mu_6$ (resp. $\mu_4$).

\uline{Supposons maintenant que} $S$ \uline{soit le spectre d'ur corps}\mbox{\begingroup\edef\encoutersep{\the\dimexpr.5\pgflinewidth\relax}\setbox0=\hbox{.}\edef\encwda{\the\wd0}\edef\enchta{\the\ht0}\edef\encdpa{\the\dp0}\rlap{\smash{\tikz[baseline=0pt]{\begin{scope}[shift={(\the\dimexpr \encwda/2\relax,\the\dimexpr -\encdpa/2+\enchta/2\relax)}]\coordinate (enc-stop-center) at (0pt,0pt);\handwritten{\draw[thin,gray,rotate around={-18:(enc-stop-center)}] (enc-stop-center) ellipse (0.57em and 1.03em);}\end{scope}\pgfresetboundingbox\path[use as bounding box] (0pt,\the\dimexpr-\encdpa\relax) rectangle (\encwda,\enchta);}}}\copy0\endgroup}

(IV) \quad \uline{Si} $3=c_4=c_6=0$, $\mathrm{Aut}(E)$ \uline{est géométriquement \unsure{isomorphe} au produit semi-direct} (\uline{non commutatif}) \uline{de} $\mathbb{Z}/(4)$ \uline{par} $\mathbb{Z}/(3)$, \uline{et l'image de} (6.6) \uline{est} $\mu_4$ .

(V) \quad \uline{Si} $2=c_4=c_6=0$, $\mathrm{Aut}(E)$ \uline{est géométriquement d'ordre} \mbox{\begingroup\edef\encoutersep{\the\dimexpr.5\pgflinewidth\relax}\setbox0=\hbox{$24$,}\edef\encwda{\the\wd0}\edef\enchta{\the\ht0}\edef\encdpa{\the\dp0}\rlap{\smash{\tikz[baseline=0pt]{\begin{scope}[shift={(\the\dimexpr \encwda/2\relax,\the\dimexpr -\encdpa/2+\enchta/2\relax)}]\coordinate (enc-order-center) at (0pt,0pt);\handwritten{\draw[thin,gray,rotate around={-38:(enc-order-center)}] ([xshift=.37em,yshift=2.08em]enc-order-center) arc[start angle=72,end angle=365,x radius=1.19em,y radius=2.2em];}\end{scope}\pgfresetboundingbox\path[use as bounding box] (0pt,\the\dimexpr-\encdpa\relax) rectangle (\encwda,\enchta);}}}\copy0\endgroup}\margin{\handwritten{\tikz[baseline=-0.3ex]{\draw[thin] (0,0.7em) -- (0.35em,-0.05em);} structure ?}} \uline{et l'image de} (6.6) \uline{est} $\mu_3$.
% archive sheet number: 29
